\section{Introduzione}

Il password spraying è uno degli attacchi più usati dagli hacker per poter aggirare i sistemi di autenticazione basati su password. Questo tipo di attacco, semplice dal punto di vista dell'esecuzione, prevede di attaccare un sistema di autenticazione sfruttando uno o più dizionari con le password più usate su tante utenze diverse\cite{pwspray}. Questo tipo di attacco è estremamente efficace, in particolare se eseguito su strutture aziendali complesse, per le quali solitamente un certo tipo di policy sulle password spinge gli utenti a sceglierle più semplici di quanto non farebbero se potessero evitare di cambiarle ogni 6 mesi .  Tendenzialmente questi attacchi di solito vengono lanciati usando botnets, ovverosia insiemi di dispositivi che sono stati compromessi che sono sfruttabili dall'attaccante per eseguire operazioni su larga scala, senza che gli utenti legittimi di tali dispositivi ne siano a conoscenza. La differenza con una tipologia di attacco a dizionario sta nella distribuzione dei tentativi, infatti in un attacco a dizionario si prova l'intero dizionario su una singola utenza per volta, utile se si vuole sbloccare un account solo , molto meno efficace se si vuole attaccare un'infrastruttura. Di fatto tutti questi attacchi sono molto pericolosi, perchè l'utente medio tende ad usare la stessa password su più utenze diverse, motivo per il quale oggi si raccomanda sempre più spesso di utilizzare password managers o sistemi di autenticazione biometrica che non soffrono di questo problema \cite{nist80063b}. L'avanzamento tecnologico rende sempre più pericolosa questa tipologia di attacchi, sopratutto perchè, anche in contesti amatoriali, è possibile trovare la tecnologia necessaria a renderli sufficientemente veloci ed efficaci.
\subsection{Analisi dei casi possibili}
I servizi di autenticazione funzionano di solito utilizzando un database, che viene interrogato mediante un servizio middleware che svolge la funzione di API di accesso, le credenziali contenute in esso vengono tipicamente criptate mediante un meccanismo di hashing, oggi il più usato è SHA256\cite{nist800107}, benchè esistano versioni più sicure come l'SHA512, queste sono meno impiegate a causa del tempo di esecuzione richiesto, che sebbene sia minimo sulla singola stringa, diventa invece oneroso quando i flussi di dati diventano importanti. Inoltre non è ancora stata trovata una collisione su SHA256 a 64 step, per cui è ancora considerato sicuro, al contrario di altre funzioni come MD5 e SHA1. I database che contengono queste chiavi, sono spesso il bersaglio di molti attaccanti, il cui obiettivo è farne un dump per poi successivamente effettuare attacchi in locale e trovare le credenziali utilizzando varie tecniche, alternativamente si può optare per attaccare direttamente il servizio di autenticazione avendo una lista di utenti che usano il servizio, in entrambi i casi il password spraying risulta essere un attacco efficace sopratutto su infrastrutture organizzative di grandi dimensioni. Supponendo che le seguenti casistiche 
\begin{itemize}
    \item Attacco tramite API usando un cluster HPC
    \item Attacco al dump usando un cluster HPC
    \item Attacco tramite API usando una botnet 
    \item Attacco al dump usando una botnet 
\end{itemize}
si può realizzare intuitivamente quali siano le architetture migliori a seconda dell'attacco. Infatti prendendo in esame il primo caso, il primo problema che si incontra è la limitazione del tasso delle richieste basate sull'indirizzo IP, inoltre con un traffico così esagerato verso un determinato bersaglio si può facilmente essere scoperti, se invece di un cluster HPC si usasse una botnet il problema verrebbe automaticamente risolto dal fatto che questa è costituita da diversi indirizzi IP per cui stategie di protezione come il rate limiter non funzionano più. Diverso è invece il caso di un attacco ad un dump, dove tipicamente il file ha grandi dimensioni e la sua trasmissione influenzerebbe in modo problematico le performance dell'attacco se si esegue tra dispositivi lontani nello spazio e con larghezze di banda molto diverse, che non è il caso però di un cluster HPC, dove molti nodi sono fortemente interconnessi tra di loro con una larghezza di banda molto più alta di quella offerta, ad esempio, dalla rete cellulare o ADSL. Per cui, usando il secondo caso come motivazione, verrà trattato proprio un attacco ad un dump in un cluster HPC, che è l'oggetto del corso.

